% Biên dịch: latexmk -lualatex -interaction=nonstopmode -halt-on-error main.tex
\documentclass[11pt,a4paper]{article}
\usepackage{fontspec}
\usepackage[main=vietnamese,english]{babel}
\usepackage{microtype}
\usepackage{amsmath}
\usepackage{libertinus}
\usepackage{booktabs,tabularx,array}
\usepackage{float}
\usepackage{tikz}
\usetikzlibrary{matrix,positioning,calc,arrows.meta}
\usepackage[most]{tcolorbox}
\usepackage{enumitem}
\usepackage{listings}
\usepackage[a4paper,margin=25mm,headheight=15pt]{geometry}
\usepackage{xcolor}
\usepackage{xurl}
\usepackage{hyperref}
\usepackage[nameinlink,noabbrev]{cleveref}
\usepackage{fancyhdr}

\definecolor{AlienNavy}{HTML}{17324D}
\definecolor{AlienTeal}{HTML}{148A8A}
\definecolor{AlienMint}{HTML}{E7F5F2}
\definecolor{AlienGold}{HTML}{D99A2B}
\definecolor{AlienRed}{HTML}{B54747}
\definecolor{AlienGray}{HTML}{F3F5F7}

\hypersetup{
  colorlinks=true,
  linkcolor=AlienNavy,
  citecolor=AlienTeal,
  urlcolor=AlienTeal,
  pdftitle={Alien Tiles - Tối thiểu hóa số lần click},
  pdfsubject={Mô tả bài toán, mô hình ILP, mô hình SAT và dữ liệu},
  pdfauthor={Tài liệu nghiên cứu prob027}
}

\pagestyle{fancy}
\fancyhf{}
\fancyhead[L]{\small\color{AlienNavy}\textsc{CSPLib prob027}}
\fancyhead[R]{\small\color{AlienNavy}Tối thiểu hóa số lần click}
\fancyfoot[C]{\small\color{AlienNavy}Trang \thepage}
\renewcommand{\headrulewidth}{0.4pt}
\renewcommand{\footrulewidth}{0pt}

\setcounter{tocdepth}{2}
\setlist[itemize]{leftmargin=*,itemsep=3pt,topsep=4pt}
\setlist[enumerate]{leftmargin=*,itemsep=5pt,topsep=4pt}
\renewcommand{\arraystretch}{1.18}
\setlength{\parindent}{0pt}
\setlength{\parskip}{5pt}

\newcommand{\Z}{\mathbb{Z}}
\newcommand{\I}{\mathcal{I}}
\newcommand{\AtMost}{\operatorname{AtMost}}
\crefname{equation}{phương trình}{các phương trình}
\crefname{table}{bảng}{các bảng}
\crefname{section}{mục}{các mục}
\crefname{subsection}{mục}{các mục}

\newtcolorbox{summarybox}[1][]{
  enhanced,breakable,colback=AlienMint,colframe=AlienTeal,
  boxrule=0.7pt,arc=2mm,left=4mm,right=4mm,top=3mm,bottom=3mm,
  fonttitle=\bfseries,title={#1}
}
\newtcolorbox{examplebox}[1][]{
  enhanced,breakable,colback=AlienGray,colframe=AlienNavy,
  boxrule=0.7pt,arc=2mm,left=4mm,right=4mm,top=3mm,bottom=3mm,
  fonttitle=\bfseries,title={#1}
}
\newtcolorbox{notebox}[1][]{
  enhanced,breakable,colback=AlienGold!8,colframe=AlienGold!85!black,
  boxrule=0.7pt,arc=2mm,left=4mm,right=4mm,top=3mm,bottom=3mm,
  fonttitle=\bfseries,title={#1}
}
\lstdefinestyle{jsonstyle}{
  basicstyle=\ttfamily\small,
  backgroundcolor=\color{AlienGray},
  frame=single,
  rulecolor=\color{AlienNavy!35},
  numbers=none,
  showstringspaces=false,
  breaklines=true,
  columns=fullflexible,
  stringstyle=\color{AlienRed}
}

\begin{document}
\hypersetup{pageanchor=false}
\begin{titlepage}
  \thispagestyle{empty}
  \begin{tikzpicture}[remember picture,overlay]
    \fill[AlienNavy] (current page.north west)
      rectangle ([yshift=-50mm]current page.north east);
    \fill[AlienTeal] ([yshift=-47mm]current page.north west)
      rectangle ([yshift=-50mm]current page.north east);
    \node[anchor=north west,text=white]
      at ([xshift=25mm,yshift=-10mm]current page.north west)
      {\Large\textsc{Tài liệu mô hình hóa}};
    \node[anchor=north west,text=white]
      at ([xshift=25mm,yshift=-22mm]current page.north west)
      {\fontsize{29}{34}\selectfont\bfseries The Alien Tiles Problem};
    \node[anchor=north west,text=white]
      at ([xshift=25mm,yshift=-39mm]current page.north west)
      {\Large CSPLib Problem 027};
  \end{tikzpicture}

  \vspace*{75mm}
  \begin{center}
  \begin{tikzpicture}[
    tile/.style={draw=AlienNavy,line width=.7pt,minimum size=9mm,
                 font=\bfseries\small,text=AlienNavy}]
    \matrix (g) [matrix of nodes,nodes={tile},
      row sep=-\pgflinewidth,column sep=-\pgflinewidth] {
      |[fill=AlienMint]| 1 & |[fill=AlienMint]| 1 & |[fill=AlienMint]| 1 & |[fill=AlienGold!25]| 0 \\
      |[fill=AlienMint]| 1 & |[fill=white]| 0 & |[fill=white]| 0 & |[fill=AlienTeal!22]| 2 \\
      |[fill=AlienGold!25]| 0 & |[fill=AlienTeal!22]| 2 & |[fill=AlienTeal!22]| 2 & |[fill=AlienTeal!22]| 2 \\
      |[fill=AlienMint]| 1 & |[fill=white]| 0 & |[fill=white]| 0 & |[fill=AlienTeal!22]| 2 \\
    };
  \end{tikzpicture}
  \end{center}

  \vfill
  \begin{tcolorbox}[
    colback=white,colframe=AlienNavy,boxrule=0pt,
    borderline west={3pt}{0pt}{AlienTeal},arc=0pt]
    \Large\bfseries Bài toán tìm lời giải có tổng số lần click nhỏ nhất\\[2mm]
    \large Ví dụ trực quan, mô hình ILP, mô hình SAT và dữ liệu
  \end{tcolorbox}
  \vspace{7mm}
  {\large Ngày cập nhật: 21 tháng 7 năm 2026\par}
\end{titlepage}
\hypersetup{pageanchor=true}

\pagenumbering{roman}
\begin{abstract}
Tài liệu chỉ xét một mục tiêu: với trạng thái đích đã cho, tìm cách click
có tổng số lần click nhỏ nhất. Trạng thái ban đầu được lấy là lưới toàn 0,
đúng theo quy ước của CSPLib prob027. Trước tiên, một ví dụ
$4\times4$ với ba màu được giải từng bước. Sau đó cùng ví dụ này được dùng
để giải thích mô hình quy hoạch tuyến tính nguyên (ILP) và mô hình SAT có
cận chi phí. Phần cuối thiết kế một quy trình sinh benchmark tái lập, mô tả
các file dữ liệu trong dự án và dẫn nguồn paper/web.
\end{abstract}

\begin{summarybox}[Bài toán trong một câu]
\textbf{Đầu vào:} kích thước $N$, số màu $c$ và target $T$.\\
\textbf{Đầu ra:} số lần click $x_{ij}$ tại mỗi ô.\\
\textbf{Mục tiêu duy nhất:}
\[
  \min \sum_{i=1}^{N}\sum_{j=1}^{N}x_{ij},
\]
với điều kiện các click phải biến lưới ban đầu toàn 0 thành đúng target
$T$.
\end{summarybox}

\begingroup
\small
\setlength{\parskip}{0pt}
\tableofcontents
\endgroup
\clearpage
\pagenumbering{arabic}

\section{Luật chơi và dữ liệu đầu vào}\label{sec:problem}

Ta có một lưới vuông $N\times N$. Mỗi ô nhận một số trong
\[
  \{0,1,\ldots,c-1\}.
\]
Các số này đại diện cho $c$ màu nằm trên một vòng tròn. Sau màu $c-1$ ta
quay lại màu 0.

Khi click ô $(i,j)$ một lần:
\begin{itemize}
  \item mọi ô trên hàng $i$ tăng thêm 1;
  \item mọi ô trên cột $j$ tăng thêm 1;
  \item ô $(i,j)$ cũng chỉ tăng 1, không phải tăng 2;
  \item tất cả phép tăng đều được lấy modulo $c$.
\end{itemize}

Ví dụ với $c=3$, màu thay đổi theo chu trình
\[
  0\longrightarrow1\longrightarrow2\longrightarrow0.
\]

Thứ tự click không quan trọng. Vì vậy, thay vì lưu một chuỗi thao tác, ta
chỉ cần lưu
\[
  x_{ij}=\text{số lần click ô }(i,j).
\]
Do click $c$ lần tại cùng một ô không tạo thay đổi ròng, chỉ cần xét
\[
  x_{ij}\in\{0,1,\ldots,c-1\}.
\]
Các tính chất này xuất phát trực tiếp từ đặc tả CSPLib
\cite{csplib027,gent2000symmetry}.

\subsection{Dạng instance dùng trong tài liệu}

Một instance gồm ba thành phần:
\[
  (N,c,T),
\]
trong đó $T=(t_{rk})$ là ma trận target. Trạng thái đầu là ma trận toàn 0.
Nếu cần trạng thái đầu $S$ khác 0, có thể thay target bằng
\[
  d_{rk}=(t_{rk}-s_{rk})\bmod c
\]
rồi giải như trường hợp bắt đầu từ 0.

\section{Ví dụ xuyên suốt: optimum bằng 3 click}\label{sec:example}

Xét lưới $4\times4$, $c=3$ và target
\begin{equation}\label{eq:example-target}
T=
\begin{bmatrix}
1&1&1&0\\
1&0&0&2\\
0&2&2&2\\
1&0&0&2
\end{bmatrix}.
\end{equation}

Một lời giải là:
\begin{itemize}
  \item click ô $(1,1)$ đúng 1 lần;
  \item click ô $(3,4)$ đúng 2 lần;
  \item không click các ô còn lại.
\end{itemize}

Ma trận click tương ứng là
\begin{equation}\label{eq:example-x}
X=
\begin{bmatrix}
1&0&0&0\\
0&0&0&0\\
0&0&0&2\\
0&0&0&0
\end{bmatrix}.
\end{equation}

\begin{examplebox}[Kiểm tra một vài ô]
\textbf{Ô $(1,2)$:} cùng hàng với $(1,1)$ nên nhận 1; không cùng hàng/cột
với $(3,4)$. Kết quả là 1.

\textbf{Ô $(2,4)$:} cùng cột với $(3,4)$ nên nhận 2. Kết quả là 2.

\textbf{Ô $(1,4)$:} nhận 1 từ $(1,1)$ và nhận 2 từ $(3,4)$:
\[
  (1+2)\bmod3=0.
\]

\textbf{Ô $(3,1)$:} nhận 2 từ $(3,4)$ và nhận 1 từ $(1,1)$:
\[
  (2+1)\bmod3=0.
\]
Các giá trị này khớp với target trong \cref{eq:example-target}.
\end{examplebox}

Tổng số click của \cref{eq:example-x} là
\[
  1+2=3.
\]
Đây chưa đủ để kết luận optimum bằng 3: ta còn phải chứng minh không có lời
giải nào dùng 0, 1 hoặc 2 click. Bộ giải SAT đi kèm đã kiểm tra các cận này:

\begin{center}
\begin{tikzpicture}[
  node distance=5mm,
  trial/.style={draw=AlienNavy,rounded corners,minimum width=28mm,
                minimum height=9mm,align=center},
  arrow/.style={-{Latex[length=2mm]},AlienTeal,thick}]
  \node[trial,fill=AlienRed!8] (k0) {$K=0$\\UNSAT};
  \node[trial,fill=AlienRed!8,right=of k0] (k1) {$K=1$\\UNSAT};
  \node[trial,fill=AlienRed!8,right=of k1] (k2) {$K=2$\\UNSAT};
  \node[trial,fill=AlienMint,right=of k2] (k3) {$K=3$\\SAT};
  \draw[arrow] (k0)--(k1);
  \draw[arrow] (k1)--(k2);
  \draw[arrow] (k2)--(k3);
\end{tikzpicture}
\end{center}

Do đó giá trị đầu tiên có lời giải là $K=3$, nên optimum bằng 3.

\begin{notebox}[Tại sao phải tối thiểu hóa?]
Một target có thể có nhiều ma trận click khác nhau. Chẳng hạn, với
$N=4,c=3$, click mỗi ô hàng 1 một lần và mỗi ô hàng 2 hai lần sẽ không làm
thay đổi lưới sau khi lấy modulo 3. Target toàn 0 vì vậy có cả lời giải
0 click và một lời giải 12 click. Hàm mục tiêu phải chọn lời giải 0 click.
\end{notebox}

\section{Mô hình toán học chung}\label{sec:math}

Đặt tập chỉ số $\I=\{1,\ldots,N\}$. Biến quyết định là
\[
  x_{ij}\in\{0,1,\ldots,c-1\},
  \qquad i,j\in\I.
\]

Xét ô target $(r,k)$. Tổng số lần ô này bị tác động là
\begin{equation}\label{eq:raw-effect}
  a_{rk}
  =
  \sum_{j\in\I}x_{rj}
  +
  \sum_{i\in\I}x_{ik}
  -
  x_{rk}.
\end{equation}

Cách đọc công thức:
\begin{enumerate}
  \item tổng thứ nhất đếm mọi click trên hàng $r$;
  \item tổng thứ hai đếm mọi click trên cột $k$;
  \item $x_{rk}$ bị đếm hai lần, nên phải trừ đi một lần.
\end{enumerate}

Để ô $(r,k)$ có đúng màu target $t_{rk}$, ta cần
\begin{equation}\label{eq:mod-constraint}
  a_{rk}\bmod c=t_{rk},
  \qquad \forall r,k\in\I.
\end{equation}

\begin{summarybox}[Mô hình tối ưu ở dạng modulo]
\begin{align}
  \min\quad
  &\sum_{i\in\I}\sum_{j\in\I}x_{ij}
  \label{eq:objective}\\
  \text{s.t.}\quad
  &\left(
    \sum_{j\in\I}x_{rj}
    +\sum_{i\in\I}x_{ik}
    -x_{rk}
  \right)\bmod c=t_{rk},
  &&\forall r,k\in\I,
  \label{eq:model-mod}\\
  &x_{ij}\in\{0,1,\ldots,c-1\},
  &&\forall i,j\in\I.
  \label{eq:model-domain}
\end{align}
\end{summarybox}

\Cref{eq:objective} là hàm mục tiêu duy nhất của tài liệu: tổng số click
càng nhỏ càng tốt. \Cref{eq:model-mod} bảo đảm lời giải tạo đúng target.

\section{Mô hình ILP}\label{sec:ilp}

ILP không dùng trực tiếp toán tử modulo. Ta thay mỗi điều kiện
\[
  a_{rk}\bmod c=t_{rk}
\]
bằng một phương trình nguyên
\begin{equation}\label{eq:quotient}
  a_{rk}-c q_{rk}=t_{rk}.
\end{equation}

$q_{rk}$ chỉ là phần thương khi chia $a_{rk}$ cho $c$; nó không phải số
lần click. Từ
\[
  0\le a_{rk}\le(2N-1)(c-1),
\]
ta có thể dùng cận an toàn
\[
  0\le q_{rk}\le
  \left\lfloor\frac{(2N-1)(c-1)}{c}\right\rfloor.
\]

\subsection{Mô hình đầy đủ}
\begin{summarybox}[ILP tối thiểu hóa số lần click]
\begin{align}
  \min\quad
  &\sum_{i\in\I}\sum_{j\in\I}x_{ij}
  \label{eq:ilp-objective}\\
  \text{s.t.}\quad
  &\sum_{j\in\I}x_{rj}
   +\sum_{i\in\I}x_{ik}
   -x_{rk}
   -c q_{rk}
   =t_{rk},
   &&\forall r,k\in\I,
  \label{eq:ilp-effect}\\
  &x_{ij}\in\Z,\quad 0\le x_{ij}\le c-1,
   &&\forall i,j\in\I,
  \label{eq:ilp-x}\\
  &q_{rk}\in\Z,\quad
   0\le q_{rk}\le
   \left\lfloor\frac{(2N-1)(c-1)}{c}\right\rfloor,
   &&\forall r,k\in\I.
  \label{eq:ilp-q}
\end{align}
\end{summarybox}

Mô hình có $N^2$ biến click $x$ và $N^2$ biến phụ $q$. Có đúng $N^2$
phương trình tác động, mỗi phương trình ứng với một ô target. Đây là dạng
tổng quát của mô hình IP tối thiểu số click được Chlond trình bày cho Alien
Tiles $7\times7$ với bốn màu \cite{chlond2006alien}.

\subsection{Đọc biến phụ q bằng ví dụ}
Quay lại target ở \cref{eq:example-target}.

\textbf{Tại ô $(1,4)$:}
\[
\underbrace{\sum_jx_{1j}}_{=1}
+
\underbrace{\sum_ix_{i4}}_{=2}
-
\underbrace{x_{14}}_{=0}
-3q_{14}
=
\underbrace{t_{14}}_{=0}.
\]
Vế tác động trước modulo là 3, nên $q_{14}=1$ và $3-3\cdot1=0$.

\textbf{Tại ô $(1,1)$:}
\[
\underbrace{\sum_jx_{1j}}_{=1}
+
\underbrace{\sum_ix_{i1}}_{=1}
-
\underbrace{x_{11}}_{=1}
-3q_{11}
=
\underbrace{t_{11}}_{=1}.
\]
Ở đây tác động đã bằng target, nên $q_{11}=0$. Hai phép tính này cho thấy
$q$ chỉ giúp ILP bỏ đi một số bội của $c$.

\section{Mô hình SAT}\label{sec:sat}

SAT cơ bản chỉ trả lời có hoặc không. Vì vậy ta đưa vào một cận $K$ và hỏi:

\begin{quote}
\emph{Có tồn tại lời giải tạo đúng target và dùng tổng cộng không quá
$K$ click hay không?}
\end{quote}

Gọi công thức này là $\Phi_K$. Giá trị tối ưu là cận $K$ nhỏ nhất làm
$\Phi_K$ khả thỏa.

\subsection{Bước 1: mã hóa số lần click ở mỗi ô}

Với mỗi ô $(i,j)$ và mỗi $v\in\{0,\ldots,c-1\}$, tạo biến Boolean
\[
  d_{ijv}=1
  \quad\Longleftrightarrow\quad
  x_{ij}=v.
\]

Ví dụ, khi $c=3$, ô $(i,j)$ có ba biến
\[
  d_{ij0},\quad d_{ij1},\quad d_{ij2}.
\]
Đúng một trong ba biến phải nhận true. Mã hóa CNF theo cặp là

\begin{align}
  &d_{ij0}\lor d_{ij1}\lor\cdots\lor d_{ij,c-1},
  \label{eq:sat-atleast}\\
  &\lnot d_{ija}\lor\lnot d_{ijb},
  \qquad 0\le a<b<c.
  \label{eq:sat-atmost}
\end{align}

\Cref{eq:sat-atleast} nói rằng phải chọn ít nhất một giá trị.
\Cref{eq:sat-atmost} cấm chọn đồng thời hai giá trị.

\subsection{Bước 2: cộng các click theo modulo}

Đối với ô target $(r,k)$, liệt kê $L=2N-1$ biến click tác động lên nó:
\[
P_{rk}=
(x_{r1},\ldots,x_{rN},
x_{1k},\ldots,x_{r-1,k},x_{r+1,k},\ldots,x_{Nk}).
\]
Biến $x_{rk}$ chỉ xuất hiện một lần trong danh sách.

Tạo biến Boolean
\[
  s_{rk,\ell,a}=1
\]
với ý nghĩa: sau khi cộng $\ell$ phần tử đầu của $P_{rk}$, phần dư hiện tại
là $a$. Tại mỗi bước $\ell$, đúng một trong
\[
  s_{rk,\ell,0},\ldots,s_{rk,\ell,c-1}
\]
được chọn.

Trạng thái đầu là phần dư 0:
\begin{equation}\label{eq:sat-start}
  s_{rk,0,0}.
\end{equation}

Giả sử phần dư cũ là $a$ và click tiếp theo có giá trị $b$. Phần dư mới
phải là $(a+b)\bmod c$. Nếu click tiếp theo là $x_{uv}$, clause chuyển là

\begin{equation}\label{eq:sat-transition}
  \lnot s_{rk,\ell-1,a}
  \lor
  \lnot d_{uvb}
  \lor
  s_{rk,\ell,(a+b)\bmod c}.
\end{equation}

Với $c=3$, quy tắc chuyển có thể đọc trực tiếp từ bảng:

\begin{table}[htbp]
\centering
\caption{Phần dư mới $(a+b)\bmod3$.}
\label{tab:mod3}
\begin{tabular}{c|ccc}
\toprule
$a\backslash b$ & 0 & 1 & 2\\
\midrule
0 & 0 & 1 & 2\\
1 & 1 & 2 & 0\\
2 & 2 & 0 & 1\\
\bottomrule
\end{tabular}
\end{table}

Sau $L$ bước, phần dư phải bằng target:

\begin{equation}\label{eq:sat-finish}
  s_{rk,L,t_{rk}},
  \qquad \forall r,k\in\I.
\end{equation}

\subsection{Bước 3: mã hóa cận tổng số click}

Ta cần chuyển giá trị $x_{ij}$ thành các đóng góp đơn vị. Với mỗi
$h=1,\ldots,c-1$, tạo

\begin{equation}\label{eq:unit-literal}
  u_{ijh}=1
  \quad\Longleftrightarrow\quad
  x_{ij}\ge h.
\end{equation}

Liên kết với các biến one-hot ở bước 1 bằng
\begin{equation}\label{eq:unit-channel}
  u_{ijh}
  \quad\Longleftrightarrow\quad
  \bigvee_{v=h}^{c-1}d_{ijv}.
\end{equation}

Với $c=3$, ta có bảng minh họa:

\begin{center}
\begin{tabular}{c|cc|c}
\toprule
$x_{ij}$ & $u_{ij1}$ & $u_{ij2}$ & Tổng đóng góp\\
\midrule
0 & 0 & 0 & 0\\
1 & 1 & 0 & 1\\
2 & 1 & 1 & 2\\
\bottomrule
\end{tabular}
\end{center}

Do đó
\begin{equation}\label{eq:unit-sum}
  \sum_{i,j}x_{ij}
  =
  \sum_{i,j}\sum_{h=1}^{c-1}u_{ijh}.
\end{equation}

Cận tổng số click không quá $K$ trở thành cardinality constraint

\begin{equation}\label{eq:sat-bound}
  \AtMost\bigl(\{u_{ijh}\},K\bigr).
\end{equation}

Ràng buộc \cref{eq:sat-bound} có thể chuyển sang CNF bằng sequential
counter \cite{sinz2005towards}. Công thức SAT đầy đủ $\Phi_K$ gồm:

\begin{enumerate}
  \item exactly-one cho các biến $d$;
  \item các trạng thái cộng modulo và clause chuyển;
  \item trạng thái cuối bằng target;
  \item cardinality constraint tổng số click không quá $K$.
\end{enumerate}

\subsection{Tìm giá trị tối ưu}

Cách dễ hiểu nhất là thử $K=0,1,2,\ldots$ cho đến khi SAT. Với ví dụ xuyên
suốt:

\[
\begin{array}{c|cccc}
K&0&1&2&3\\
\hline
\Phi_K&\text{UNSAT}&\text{UNSAT}&\text{UNSAT}&\text{SAT}
\end{array}
\]

Vì $K=3$ là cận nhỏ nhất khả thỏa, lời giải có tổng click bằng 3 là tối ưu.
Khi miền cận lớn, có thể dùng tìm kiếm nhị phân hoặc giữ một solver SAT
incremental để tái sử dụng các clause đã học.

\section{Dataset và nguồn tham khảo}\label{sec:data}

\subsection{Nguồn chính}

\begin{itemize}
  \item \textbf{CSPLib prob027} cung cấp luật chơi và các mô hình Essence
        \cite{csplib027,csplib027models}.
  \item \textbf{Chlond (2006)} trình bày trực tiếp mô hình IP tối thiểu hóa
        số click cho Alien Tiles $7\times7$, bốn màu
        \cite{chlond2006alien}.
  \item \textbf{Gent, Linton và Smith (2000)} công bố một ma trận click
        $4\times4$, $c=3$ có tổng 10; target suy ra từ ma trận này được dùng
        như một instance cố định trong bộ dữ liệu
        \cite{gent2000symmetry,csplib027results}.
\end{itemize}

Trang \emph{Data files} của CSPLib hiện không liệt kê file instance nào
\cite{csplib027data}. Vì vậy các file JSON trong dự án là bộ dữ liệu thực
nghiệm đi kèm, không phải bộ dữ liệu chính thức của CSPLib.

\subsection{Mục tiêu thiết kế benchmark}

Một bộ benchmark dùng cho khóa luận cần đáp ứng đồng thời bốn yêu cầu:

\begin{enumerate}
  \item có instance chắc chắn có lời giải để đo khả năng tối ưu hóa;
  \item có target không bị ép phải có lời giải, nhờ đó vẫn đánh giá được
        khả năng phát hiện UNSAT;
  \item phủ nhiều giá trị $N$, $c$ và nhiều mức độ thưa/dày;
  \item có thể sinh lại đúng từng byte từ cấu hình và seed đã công bố.
\end{enumerate}

\begin{notebox}[Điểm dễ gắn nhãn sai]
Nếu target được sinh từ ma trận click $X^{\mathrm w}$ có tổng $U$, ta chỉ
biết optimum không vượt quá $U$. Có thể tồn tại ma trận click khác có tổng
nhỏ hơn. Vì vậy generator ghi \texttt{known\_upper\_bound = U}, không ghi
\texttt{known\_optimum = U}.
\end{notebox}

\subsection{Hai họ instance bổ sung nhau}

\subsubsection{Họ witness: có lời giải theo cách xây dựng}

Với mỗi bộ tham số $(N,c,\rho)$, đặt
\[
  s=\max\!\left(1,\operatorname{round}(\rho N^2)\right).
\]
Generator chọn đều $s$ ô khác nhau. Tại mỗi ô được chọn, số click được lấy
đều trong $\{1,\ldots,c-1\}$; các ô còn lại bằng 0. Ma trận nhận được gọi
là witness $X^{\mathrm w}$. Target được tính theo đúng luật chơi:
\begin{equation}\label{eq:data-witness}
  T_{rk}=\left(
    \sum_j x^{\mathrm w}_{rj}
    +\sum_i x^{\mathrm w}_{ik}
    -x^{\mathrm w}_{rk}
  \right)\bmod c.
\end{equation}

Instance này chắc chắn khả thi vì $X^{\mathrm w}$ tạo ra $T$. Tham số
$\rho$ điều khiển tỷ lệ ô click khác 0 trong witness, chứ không trực tiếp
điều khiển độ thưa của target.

\subsubsection{Họ uniform-target: chưa biết SAT hay UNSAT}

Mỗi phần tử target được lấy độc lập và đều:
\begin{equation}\label{eq:data-uniform}
  T_{ij}\sim\operatorname{Uniform}\{0,1,\ldots,c-1\}.
\end{equation}
Generator không tạo witness và gắn trạng thái
\texttt{unclassified}. Sau đó solver chính xác mới phân loại instance thành
OPTIMAL, UNSAT hoặc TIMEOUT.

\begin{table}[H]
\centering
\caption{Vai trò của hai họ instance.}
\label{tab:data-families}
\begin{tabularx}{\textwidth}{@{}l X X@{}}
\toprule
\textbf{Họ} & \textbf{Biết ngay khi sinh} & \textbf{Dùng để đánh giá} \\
\midrule
Witness
& Có lời giải; tổng click của witness là upper bound
& Khả năng tìm và chứng minh nghiệm tối ưu trên các mật độ click khác nhau \\
Uniform-target
& Chỉ biết target hợp lệ về định dạng
& Cả tối ưu hóa trên instance khả thi và phát hiện instance không có lời giải \\
\bottomrule
\end{tabularx}
\end{table}

Chỉ dùng họ witness sẽ làm bộ dữ liệu lệch hoàn toàn về phía instance khả
thi. Chỉ dùng họ uniform-target lại không bảo đảm có đủ instance khả thi ở
từng cấu hình. Hai họ vì vậy phải được giữ riêng trong báo cáo kết quả.

\subsection{Lấy mẫu theo tầng và chống thiên lệch}

Mỗi cấu hình là một tầng độc lập. Cấu hình benchmark khuyến nghị ban đầu là:

\begin{table}[H]
\centering
\caption{Cấu hình benchmark khuyến nghị cho thực nghiệm chính.}
\label{tab:data-config}
\begin{tabularx}{\textwidth}{@{}l X l@{}}
\toprule
\textbf{Thành phần} & \textbf{Giá trị} & \textbf{Số mức} \\
\midrule
Kích thước $N$ & $4,5,6,8,10,12$ & 6 \\
Số màu $c$ & $2,3,4$ & 3 \\
Mật độ witness $\rho$ & $0.10,0.30,0.60$ & 3 \\
Họ không có $\rho$ & uniform-target & 1 \\
Số lặp mỗi tầng & 2 pilot và 8 evaluation & 10 \\
\bottomrule
\end{tabularx}
\end{table}

Với mỗi cặp $(N,c)$ có ba tầng witness và một tầng uniform-target. Tổng số
instance là
\[
  6\times3\times(3+1)\times(2+8)=720,
\]
gồm 144 instance pilot và 576 instance evaluation. Các instance $N=3$ chỉ
dùng trong bộ kiểm chứng nhỏ; không đưa vào benchmark chính vì với $c=2$
không gian target khả thi rất nhỏ, dễ lặp lại target giữa các tầng.
Bộ pilot được phép dùng
để sửa lỗi, chọn encoding và tham số solver. Bộ evaluation chỉ được chạy
sau khi cấu hình thuật toán đã chốt; đây là số liệu dùng cho kết luận cuối.

Để việc sinh dữ liệu không phụ thuộc thứ tự vòng lặp, mỗi instance dùng một
seed 64-bit riêng:
\[
  \operatorname{seed}_{\I}
  =\operatorname{first64}\!\left(
  \operatorname{SHA256}(S,\text{split},N,c,\text{họ},\rho,r,a)\right),
\]
trong đó $S$ là master seed, $r$ là số thứ tự lặp và $a$ là số lần thử lại
khi gặp target trùng. Target được nhận diện bằng SHA-256 của bộ $(N,c,T)$.
Nếu hash đã xuất hiện trong bộ dữ liệu thì generator tăng $a$ và sinh lại.
Các phép lấy mẫu tiếp theo dùng bộ sinh counter-mode SHA-256 do dự án cài
đặt, vì vậy kết quả không phụ thuộc vào thay đổi nội bộ của
\texttt{random.Random} giữa các phiên bản Python.

\subsection{Nhãn và quy trình sinh}

Mỗi instance đi qua các bước sau:

\begin{enumerate}
  \item đọc một tầng $(N,c,\text{họ},\rho,\text{split},r)$;
  \item dẫn xuất instance seed từ master seed;
  \item sinh witness hoặc sinh trực tiếp target;
  \item tính target, thống kê tần suất màu và hash;
  \item loại target trùng trong cùng bộ benchmark;
  \item kiểm tra kích thước, miền giá trị và kiểm tra lại witness;
  \item ghi JSON, \texttt{manifest.csv} và
        \texttt{dataset\_summary.json}.
\end{enumerate}

Quy tắc gắn nhãn được cố định như sau:

\begin{itemize}
  \item witness: \texttt{feasible\_by\_construction} và có
        \texttt{known\_upper\_bound};
  \item uniform-target: \texttt{unclassified};
  \item solver tìm được lời giải $K$ click nhưng chưa chứng minh tối ưu:
        chỉ cập nhật upper bound;
  \item solver chứng minh tối ưu: mới ghi \texttt{known\_optimum = K};
  \item solver chứng minh vô nghiệm: ghi \texttt{known\_status = UNSAT};
  \item hết thời gian: ghi TIMEOUT, không được đổi thành UNSAT.
\end{itemize}

\subsection{Triển khai trong dự án}

Generator chỉ dùng thư viện chuẩn Python và từ chối ghi đè thư mục đã tồn
tại. Lệnh sinh bộ demo đi kèm là:

\begin{lstlisting}[style=jsonstyle,language={}]
python3 experiments/generate_instances.py \
  --output-dir data/generated_demo \
  --sizes 3 4 --colours 2 3 \
  --densities 0.15 0.40 \
  --pilot-count 1 --evaluation-count 2 \
  --master-seed 270027
\end{lstlisting}

Bộ demo có 36 instance: 12 pilot, 24 evaluation; trong đó 24 instance
witness và 12 instance uniform-target. Tất cả 36 target là khác nhau.
Kiểm tra toàn bộ dataset bằng:

\begin{lstlisting}[style=jsonstyle,language={}]
python3 experiments/validate_dataset.py data/generated_demo
\end{lstlisting}

Lệnh sinh bộ 720 instance khuyến nghị là:

\begin{lstlisting}[style=jsonstyle,language={}]
python3 experiments/generate_instances.py \
  --output-dir data/benchmark_v1 \
  --sizes 4 5 6 8 10 12 --colours 2 3 4 \
  --densities 0.10 0.30 0.60 \
  --pilot-count 2 --evaluation-count 8 \
  --master-seed 270027
\end{lstlisting}

Các file được tổ chức theo
\texttt{split/family/nNN\_cCC/}. File \texttt{manifest.csv} là chỉ mục dùng
cho chương trình chạy thực nghiệm; \texttt{dataset\_summary.json} lưu toàn
bộ cấu hình để có thể tái lập.

\subsection{Các instance kiểm chứng thủ công}

\begin{table}[H]
\centering
\caption{Bộ dữ liệu dùng để thử mô hình tối thiểu số click.}
\label{tab:dataset}
\footnotesize
\begin{tabularx}{\textwidth}{@{}>{\ttfamily}l c c X l@{}}
\toprule
\normalfont\textbf{File} & $N$ & $c$ &
\normalfont\textbf{Mô tả} & \normalfont\textbf{Kết quả} \\
\midrule
3x3\_c2\_allones.json
&3&2&Target toàn 1&minimum 3\\
4x4\_c3\_allones.json
&4&3&Target toàn 1&minimum 4\\
4x4\_c3\_easy.json
&4&3&Ví dụ xuyên suốt ở \cref{sec:example}&minimum 3\\
4x4\_c3\_infeasible.json
&4&3&Instance âm để kiểm tra solver&không có lời giải\\
4x4\_c3\_hardest\_gent2000.json
&4&3&Target suy ra từ witness đã công bố&minimum 10
\cite{gent2000symmetry}\\
7x7\_c4\_solvable.json
&7&4&Target tự sinh từ witness 6 click&upper bound 6\\
\bottomrule
\end{tabularx}
\end{table}

Các giá trị minimum 3, 4 và 10 đã được kiểm tra bằng
\texttt{models/sat\_variant2.py}. Với file $7\times7$, dự án mới gắn nhãn
upper bound 6; không nên gọi 6 là optimum nếu chưa chứng minh không tồn tại
lời giải nhỏ hơn.

\subsection{Schema dữ liệu sinh tự động}

Một instance witness sinh tự động có các trường chính sau:

\begin{lstlisting}[style=jsonstyle,language={},caption={Schema JSON rút gọn của một instance witness.}]
{
  "schema_version": "1.0",
  "name": "evaluation_n04_c03_witness_d150_r000",
  "N": 4,
  "c": 3,
  "target": [[...], [...], [...], [...]],
  "generation": {
    "method": "witness",
    "split": "evaluation",
    "master_seed": 270027,
    "instance_seed": 123456789,
    "target_sha256": "..."
  },
  "guaranteed_solvable": true,
  "status_at_generation": "feasible_by_construction",
  "witness_clicks": [[...], [...], [...], [...]],
  "witness_total_clicks": 5,
  "known_upper_bound": 5
}
\end{lstlisting}

Instance uniform-target không có ba trường witness và upper bound ở cuối.
Trường \texttt{known\_optimum} cũng không xuất hiện khi sinh; nó chỉ được
bổ sung sau khi một solver chính xác đã chứng minh giá trị tối ưu.

\clearpage
\section{Cách chạy mô hình SAT đi kèm}\label{sec:run}

Script tối ưu target cố định là
\texttt{models/sat\_variant2.py}. Chạy ví dụ xuyên suốt bằng:

\begin{lstlisting}[style=jsonstyle,language={}]
python3 models/sat_variant2.py \
  --input data/4x4_c3_easy.json
\end{lstlisting}

Tham số \texttt{--input-dir} đọc đệ quy các file instance và tự bỏ qua
\texttt{dataset\_summary.json}. Ví dụ chạy toàn bộ tập pilot đã sinh:

\begin{lstlisting}[style=jsonstyle,language={}]
python3 models/sat_variant2.py \
  --input-dir data/generated_demo/pilot
\end{lstlisting}

Kết quả cần kiểm tra gồm:

\begin{itemize}
  \item \texttt{Optimal X}: ma trận số lần click;
  \item \texttt{Minimum total clicks}: giá trị hàm mục tiêu;
  \item \texttt{Verification: PASSED}: áp dụng lại ma trận click tạo đúng
        target.
\end{itemize}

Với instance không có lời giải, script trả về
\texttt{UNSATISFIABLE}; khi đó không tồn tại giá trị minimum.

\section{Kết luận}\label{sec:conclusion}

Bài toán được nghiên cứu trong tài liệu có một mục tiêu rõ ràng:
\[
  \min\sum_{i,j}x_{ij}.
\]
Các ràng buộc hàng, cột và modulo bảo đảm ma trận click tạo đúng target.
Trong ILP, modulo được thay bằng biến thương nguyên $q_{rk}$. Trong SAT,
ta hỏi lần lượt liệu có lời giải với tổng click không quá $K$ hay không;
cận $K$ nhỏ nhất trả về SAT chính là optimum.

Ví dụ $4\times4$, $c=3$ cho thấy toàn bộ quá trình: ma trận click có tổng
3 tạo đúng target, còn các cận 0, 1 và 2 đều UNSAT. Vì vậy minimum bằng 3.

\clearpage
\bibliographystyle{plainurl}
\bibliography{references/references}

\end{document}
